% !TEX root = ../Hauptdatei.tex
% ============================================================
% Kapitel: Grundlegende Tabellen
% ============================================================

\chapter{Grundlegende Tabellen}
\label{sec:grundlagen}
\index{Tabelle!Grundlagen}

\section{Einfache Tabelle}

Eine minimale Tabelle mit zentrierten Spalten:

\begin{table}[htbp]
    \centering
    \caption{Eine einfache Tabelle}
    \label{tab:einfach}
    \begin{tabular}{ccc}
        \hline
        Name & Alter & Stadt \\
        \hline
        Anna & 25 & Berlin \\
        Ben & 30 & Hamburg \\
        Clara & 28 & München \\
        \hline
    \end{tabular}
\end{table}

\section{Spaltenausrichtung}

Die drei Standard-Ausrichtungen\index{Tabelle!Spaltenausrichtung}:

\begin{table}[htbp]
    \centering
    \caption{Spaltenausrichtungen: links, zentriert, rechts}
    \label{tab:ausrichtung}
    \begin{tabular}{lcr}
        \hline
        Linksbündig & Zentriert & Rechtsbündig \\
        \hline
        Text & Mitte & 42 \\
        Noch mehr & Hier & 7 \\
        \hline
    \end{tabular}
\end{table}

% ============================================================
\chapter{Professionelle Tabellen mit booktabs}
\label{sec:booktabs}
\index{booktabs (Paket)}
% ============================================================

Das Paket \texttt{booktabs}\cite{mittelbach2004} bietet horizontale Linien in verschiedenen Dicken:

\begin{table}[htbp]
    \centering
    \caption{Professionelle Tabelle mit \texttt{booktabs}}
    \label{tab:booktabs}
    \begin{tabular}{lrr}
        \toprule
        Fachsemester & Studierende & Anteil (\%) \\
        \midrule
        1.\ Semester & 432 & 24{,}0 \\
        2.\ Semester & 387 & 21{,}5 \\
        3.\ Semester & 341 & 18{,}9 \\
        4.\ Semester & 298 & 16{,}6 \\
        Höher & 344 & 19{,}1 \\
        \midrule
        Gesamt & 1802 & 100{,}0 \\
        \bottomrule
    \end{tabular}
\end{table}

\section{Teilstriche mit booktabs}

\begin{table}[htbp]
    \centering
    \caption{Tabelle mit \texttt{cmidrule}\index{cmidrule}}
    \label{tab:cmidrule}
    \begin{tabular}{llrr}
        \toprule
        Kategorie & Fach & Teilnehmer & Bestanden (\%) \\
        \midrule
        \multirow{3}{*}{Pflicht}
            & Mathematik & 120 & 78 \\
            & Physik & 95 & 82 \\
            & Informatik & 110 & 91 \\
        \cmidrule(lr){2-4}
        \multirow{2}{*}{Wahl}
            & Philosophie & 34 & 88 \\
            & Kunst & 28 & 95 \\
        \bottomrule
    \end{tabular}
\end{table}

% ============================================================
\chapter{Zellen verbinden}
\label{sec:verbinden}
\index{Tabelle!Zellen verbinden}
% ============================================================

\section{Spalten verbinden (multicolumn)}

\begin{table}[htbp]
    \centering
    \caption{Spalten mit \texttt{multicolumn}\index{multicolumn} verbunden}
    \label{tab:multicolumn}
    \begin{tabular}{llll}
        \toprule
        \multicolumn{2}{c}{Name} & Vorname & Note \\
        \midrule
        Müller & Anna & Informatik & 1{,}3 \\
        Schmidt & Ben & Physik & 2{,}0 \\
        Weber & Clara & Mathematik & 1{,}7 \\
        \bottomrule
    \end{tabular}
\end{table}

\section{Zeilen verbinden (multirow)}

\begin{table}[htbp]
    \centering
    \caption{Zeilen mit \texttt{multirow}\index{multirow} verbunden}
    \label{tab:multirow}
    \begin{tabular}{llr}
        \toprule
        Gruppe & Aufgabe & Punkte \\
        \midrule
        \multirow{3}{*}{Gruppe A}
            & Aufgabe 1 & 15 \\
            & Aufgabe 2 & 18 \\
            & Aufgabe 3 & 12 \\
        \midrule
        \multirow{3}{*}{Gruppe B}
            & Aufgabe 1 & 20 \\
            & Aufgabe 2 & 14 \\
            & Aufgabe 3 & 16 \\
        \bottomrule
    \end{tabular}
\end{table}

\section{Zeilen und Spalten gleichzeitig}

\begin{table}[htbp]
    \centering
    \caption{Kombination aus \texttt{multirow} und \texttt{multicolumn}}
    \label{tab:kombiniert}
    \begin{tabular}{llcc}
        \toprule
        & & \multicolumn{2}{c}{Prüfung} \\
        \cmidrule(lr){3-4}
        Fach & Dozent & Klausur & Mündlich \\
        \midrule
        \multirow{2}{*}{Mathematik}
            & Prof.\ A & 1{,}7 & 2{,}0 \\
            & Prof.\ B & 2{,}3 & 1{,}9 \\
        \midrule
        Physik & Prof.\ C & 2{,}0 & 1{,}5 \\
        \bottomrule
    \end{tabular}
\end{table}

% ============================================================
\chapter{Farbige Tabellen}
\label{sec:tab-farbig}
\index{Tabelle!Farbig}
% ============================================================

\section{Farbige Zeilen}

\begin{table}[htbp]
    \centering
    \caption{Tabelle mit farbigen Zeilen}
    \label{tab:farbige-zeilen}
    \begin{tabular}{lrr}
        \toprule
        \rowcolor{blue!15}
        Semester & Studierende & Anteil (\%) \\
        \midrule
        WiSe 24/25 & 420 & 23{,}3 \\
        \rowcolor{gray!10}
        SoSe 25 & 380 & 21{,}1 \\
        WiSe 25/26 & 410 & 22{,}8 \\
        \rowcolor{gray!10}
        SoSe 26 & 350 & 19{,}4 \\
        \midrule
        \rowcolor{blue!8}
        Gesamt & 1560 & 86{,}6 \\
        \bottomrule
    \end{tabular}
\end{table}

\section{Farbige Spalten}

\begin{table}[htbp]
    \centering
    \caption{Tabelle mit farbiger Spalte}
    \label{tab:farbige-spalte}
    \begin{tabular}{l >{\columncolor{green!10}}r r}
        \toprule
        Fach & Punkte & Bestanden \\
        \midrule
        Mathe & 85 & Ja \\
        Physik & 72 & Ja \\
        Chemie & 48 & Nein \\
        \bottomrule
    \end{tabular}
\end{table}

% ============================================================
\chapter{Spaltenbreite und Textumbruch}
\label{sec:spaltenbreite}
\index{Tabelle!Spaltenbreite}
% ============================================================

\section{Feste Spaltenbreite mit Umbruch}

Mit \texttt{p\{...\}}\index{p-Spalte} wird der Text in der Spalte umgebrochen:

\begin{table}[htbp]
    \centering
    \caption{Tabelle mit fester Spaltenbreite}
    \label{tab:breite}
    \begin{tabular}{p{3cm} p{5cm} r}
        \toprule
        Fach & Beschreibung & ECTS \\
        \midrule
        Mathematik I & Analysis, Lineare Algebra und numerische Methoden & 8 \\
        Physik & Mechanik, Thermodynamik und Elektrodynamik & 6 \\
        Informatik & Programmierung, Algorithmen und Datenstrukturen & 7 \\
        \bottomrule
    \end{tabular}
\end{table}

\section{Zentrierter Umbruch mit m-Spalten}

\begin{table}[htbp]
    \centering
    \caption{Vertikal zentrierte \texttt{m}-Spalten\index{m-Spalte}}
    \label{tab:m-spalten}
    \begin{tabular}{m{3cm} m{5cm} r}
        \toprule
        Fach & Beschreibung & ECTS \\
        \midrule
        Mathe & Lange Beschreibung die umbricht & 8 \\
        \bottomrule
    \end{tabular}
\end{table}

% ============================================================
\chapter{Lange Tabellen}
\label{sec:longtable}
\index{longtable (Paket)}
% ============================================================

Für Tabellen, die über mehrere Seiten gehen, eignet sich \texttt{longtable}:

\begin{longtable}{lrr}
    \toprule
    Fach & Teilnehmer & Bestanden (\%) \\
    \midrule
    \endfirsthead

    \toprule
    Fach & Teilnehmer & Bestanden (\%) \\
    \midrule
    \endhead

    \midrule
    \multicolumn{3}{r}{\textit{Fortsetzung auf der nächsten Seite}} \\
    \endfoot

    \bottomrule
    \endlastfoot

    Mathematik I & 120 & 78 \\
    Mathematik II & 95 & 72 \\
    Physik I & 110 & 80 \\
    Physik II & 85 & 75 \\
    Informatik I & 150 & 85 \\
    Informatik II & 130 & 82 \\
    Chemie & 70 & 70 \\
    Biologie & 65 & 88 \\
    Philosophie & 40 & 92 \\
    Kunstgeschichte & 35 & 95 \\
    Geschichte & 45 & 90 \\
    Soziologie & 55 & 87 \\
    Psychologie & 80 & 83 \\
    Wirtschaft & 100 & 79 \\
    Recht & 90 & 76 \\
\end{longtable}